\documentclass[a4paper]{book}
\usepackage{listings}
\usepackage[margin=50pt]{geometry}
\usepackage{layout}
\usepackage{enumitem}
\usepackage{makecell}
\usepackage[T1]{fontenc}
\usepackage{hyperref}
\newcommand\tab[1][1cm]{\hspace*{#1}}
\title{Endworld PnP}
\author{Maric}
\date{}
\setcounter{secnumdepth}{4}
\setcounter{tocdepth}{4}
\hypersetup{
    colorlinks=true,
    linkcolor=blue,
    filecolor=magenta,
    urlcolor=cyan,
}
\urlstyle{same}
\raggedbottom
\begin{document}
\maketitle
\tableofcontents
\chapter{The EndWorld}\label{ch:theendworld}
Years have gone by since the Great War.
Many Years.
Even though few remember how long it has been, the remnants of a time long gone still remind us that once,
long ago, the world has been a nicer place.
But, the world we know now \ldots is not a nice place.
Such a shame.
I have lived a long life, and I have learned to adapt to the new World,
but just as the Contamination eats away at technology and the lives of those who come into contact with it, time
eats away at me.
Such a shame, there is so little of it left for me.
But the plans are in motion, immortality might be in my grasp.\par
I am \ldots well, it seems some Details are starting to slip my mind.
I am Quaesitor, and yet again I am wrong, but not entirely.
I was Quaesitor.
I traveled from the great city of Terracina to the quaint outpost of Kodal to learn.
A lifetime of learning of trying to protect what the Supervivo have protected so long.
The ancient Knowledge and ancient Arts.
Microelectronics.
Quantum Physics.
The like.
We, the Survivors, the Triumphant, the Conquerors,
forced to huddle in what were mere shelters for a storm that has not subsided over the centuries,
it is such a shame.
Marvels of modern technology have kept me from dying these past years since I retired,
but no army of high-tech mechs, no experimental biotech treatment seems to be able to extend my mind.
I have forgotten so much.
\par Such a shame.\par
Of course we are not the only ones who survived.
Beasts stalk the lands, ferocious beasts, Megafauna, sly, intelligent and hungry creatures.
In some cases literally.
But even Humanoids, a pale imitation of real Humans, the descendants of vile mutants,
military personnel enhancement projects, the failed attempts of foolish hope,
are vile mutants and savages themselves, such a shame.
They raid caravans, keep to their disjointed barbaric cultures,
forever cursed to harbor the very Contamination killing their chances of ever knowing the pleasures of technology.
And it seems that they have adapted to the Contamination, so it does not even kill them as much these days.
Such a shame. \par
Alacast can purify them just like anyone or anything else,
it changes the Contamination back to more normal dangers, but dangers that can be dealt with by proper humans.
Tribals can live without the Contamination,
if ever someone would invest in cleaning a Savage up that much, they likely still can not contribute.
Such a shame.
I have tried, yet they remained uncooperative, another sign of their inferiority.
Some babble about cages and insults, barely comprehensible ramblings in their animalistic mockery of
speech.
They do not usually last long once they have gone through decon.
Such a shame.
Especially those infuriating individuals who seem to have allied with the Contamination trying to kill us all,
wielding its Zones Of Normalcy Exclusion like weapons.
Such a shame.
Wide ZONEs of heightened gravity or brittleness, Oxygen free ZONEs or pure Oxidation.
Such a shame that they have to appear at all even if caused by natural sources.
Savages funneling Contamination into a ZONE  until even the best seal breaks down and the Pilots are eaten alive
\ldots \par Such a Shame. \par \par
-- Transcript of the only Conversation the oldest Document in the Great Library of Kodal - a severely damaged
upload of the Brain of an ancient Scientist - will have.\newpage


\chapter{Character Creation}\label{ch:charactercreation}
First, get a general idea who your character is and what they do.
Note your Character name, your name, the name of the
Campaign and your concept for the character on the character sheet.
This will help you choose in the following steps.

\section{Step 1 - Categories}\label{sec:step1-categories}

Choose Category, Fate and Affinity\par
\begin{enumerate}
    \item \textbf{Resistance} describes how able the Character is to survive the hardships of the outside world
    \item \textbf{Fate} is ``Beginners Luck'' and the room to grow of the character, separating NPC from PC\@.
    \item \textbf{Affinity} for technology allows usage of complex and futuristic gadgets
\end{enumerate}
all start at 1, then 7 points are divided between them, maximum per is 5.

\section{Step 2 - Scores}\label{sec:step2-scores}
\subsection{Abilities}\label{subsec:abilities}
Order Social, Physical, Mental, Spiritual and Ability in the order that you think they are most
usefull to your character concept. \par
In your primary priority you get 16 points, in your secondary priority you get 12 points and in your tertiary
priority you get 8 points.
The remaining two Priorities get 4 points.\par
Skills range from 1--3:
\begin{enumerate}
    \item proficiency
    \item mastery
    \item special expertise
\end{enumerate}

\subsection{Attributes}\label{subsec:attributes}
Attributes, which range from 1--5 can be unlocked by allocating spent skill levels.\par
To raise an attribute a level, 1 point per current level is required, so the levels have requirements of 1;2;3;4;5.
With total costs of the levels being:\par
\begin{enumerate}
    \item 1
    \item 3
    \item 6
    \item 10
    \item 15
\end{enumerate}
To have an Attribute at Level 5, 15 total levels in Skills (8 $\cdot$ 2 or 15 $\cdot$ 1 or anything between) are
required.
To reach Attributes of 5/5/5 120 total levels of skills would be required.
Attributes can be 0, making the character unable to attempt rolls that would use that attribute.

\subsection{Finalizing}\label{subsec:finalizing}
All skillpoints need to be spent, unallocated skillpoints remain to be allocated to Attributes later.
Reskilling / attributing a new character can be done if player and Storyteller agree.
Every level 1 skill requires some sort of mention in the backstory of the character, level 2 skills should be given
an explicit reason and level 3 skills should have an explanation of how the character has reached peak human ability
in that skill.

Later Points are awarded to specific Skills through roleplay and spending Off-time training.
These Skillpoints can then be allocated to an Attribute at any time.

\section{Step 3 - Finishing touches}\label{sec:step3-finishingtouches}
You get 100 Credits, you can spend these on:
\begin{enumerate}[label = -]
    \item Your \hyperref[ch:mechs]{Mech}
    \item \hyperref[sec:tradegoods]{Goods} (3 per credit)
    \item Gear, Spells, etc.
    \item \hyperref[sec:step2-scores]{Bonus Character creation points} (5:1)
\end{enumerate}


\section{Rules}\label{sec:rules}
\chapter{Rules}\label{ch:rules}
\section{The Selector System}\label{sec:selectorsystem}
This system uses ten sided dice.
A \emph{Check} is a roll of 5 dice interpreted in a certain way,
with the result describing how successfull the attempted action was.\par
Whenever there is advantage or disadvantage of any kind,
\hyperref[sec:bonus--and-penaltydice]{Bonus- or Penaltydice} may be applicable,
increasing the number of dice rolled.
The roll itself is interpreted by your \emph{Selectors} which usually are defined
by the applicable attribute and skills, but may be static or equipment based.
The dice in the roll are ordered in ascending order.
To get the result of the roll (or to \emph{interpret} the roll), every selector \emph{selects} describes
the number of the die (from the bottom).

Selected dice are then added.

Selectors greater than 5 simply select the 5th die.
A selector of 0 or less selects no die.
The usual number of selectors is 2.\par
The effects and result of the Check may then be further modified by \hyperref[sec:resonance]{Resonance}.

\section{Bonus- and Penaltydice}\label{sec:bonus--and-penaltydice}
Bonus- and Penaltydice cancel each other out before the roll.\par
Any remaining Bonus- or Penaltydice are rolled with the normal dice.
If an amount of Bonusdice were rolled,
the same amount of the lowest dice of the roll are removed.
For Penaltydice, the highest are removed instead.
Neither Bonus nor Penaltydice ever change the number of dice being interpreted, just the relative Chances.

\section{Resonance}\label{sec:resonance}
Resonance is when in a Roll, more than one die shows the same Number.\par
The Resonance \emph{amplitude} is the number of dice \emph{that exceed 1}, meaning it is lower by 1 than the total
number of dice.\par
The Resonance \emph{frequency} is the number those dice show.
If no effect states anything about a Resonance, they have no effect.

\chapter{Skills and Attributes}\label{ch:skillsandattributes}
\section{Physical Skills}\label{sec:strength-skills}
\subsection{Heavy Athletics}\label{subsec:heavy-athletics}
do you even lift?
\subsection{Running}\label{subsec:running}
running fast is usually fitness, but agility may come in if it is running to avoid being shot at
\subsection{Jumping}\label{subsec:jumping}
leaping high or far, this one has to do with energy bursts from the legs
\subsection{Throwing}\label{subsec:throwing}
yeet
\subsection{Survival}\label{subsec:survival}
wilderness skills
\subsection{Acrobatics}\label{subsec:acrobatics}
Uses include: doging from cover to cover, backflips, climbing, balancing and many others
\subsection{Weapon()}\label{subsec:weapon}
There are many weapons to choose from!
\subsection{Footwork}\label{subsec:footwork}
Used mainly in close range, \textbf{footwork} describes the characters ability to move fast and decisively in combat,
outmaneuvering would be enemies

\section{Social Skills}\label{sec:social-skills}
\subsection{Trade}\label{subsec:trade}
trading is used to \hyperref[ch:trade]{Trade}
\subsection{Diplomacy}\label{subsec:diplomacy}
finding solutions that both sides want
\subsection{Misdirection}\label{subsec:misdirection}
includes lying and general providing of wrong information as correct
\subsection{Style}\label{subsec:style}
panache!
style!
generally being impressive/following characters style
\subsection{Intimidation}\label{subsec:intimidation}
\subsection{Rhethoric}\label{subsec:rhethoric}
Words, I have the best Words
\subsection{Command}\label{subsec:command}
giving orders to others
\subsection{Empathy}\label{subsec:empathy}
knowing how others feel
\subsection{Etiquette}\label{subsec:etiquette}
Knowing how to behave
\subsection{Rumor}\label{subsec:rumor}
the skill at navigating the rumor mill, sifting through verbally repeated information and even starting rumors
\subsection{Art()}\label{subsec:art}
may be any art that is used primarily for self expression and social interaction.

\section{Mental Skills}\label{sec:wisdom-skills}
\subsection{Healing}\label{subsec:healing-skill}
important skill for \hyperref[subsec:treatment]{Treating of Wounds}
\subsection{Research}\label{subsec:research}
acquiring and processing information from readily available sources
\subsection{Instinct}\label{subsec:instinct}
used for sensing things that are not being paid attention to and as situational awareness in combat
\subsection{Search}\label{subsec:search}
user for sensing things specifically declared by the player or looked for
\subsection{Strategy}\label{subsec:strategy}
larger ``zoomed out'' decisions and their repercussions
\subsection{Tactics}\label{subsec:tactics}
local ``zoomed in'' decisions and their repercussions
also sometimes used with agility for advancing / retreating in a firefight
\subsection{Navigation}\label{subsec:navigation}
if you need to know where you are and how to get to a place
\subsection{Knowledge()}\label{subsec:knowledge}
lore or science, any sort of mostly theoretical knowledge

\section{Ability Skills}\label{sec:ability-skills}
\subsection{Biotech}\label{subsec:biotech}
Biotech is the Fusion of Biology and Technology \par
The ancient texts define Biotechnology as "any technological application that uses biological systems, living organisms,
or derivatives thereof, to make or modify products or processes for specific use".
Biotechnology is such a diverse group of skills, that each field is ungrouped rules-wise.
\begin{enumerate}[label= -]
    \item \textbf{Red} Biotechnology Encompasses genetics, genomics, pharmacology, but also medical advances like
    genetherapy and automedics and deconamination of living things.
    \item \textbf{Green} Biotechnology is used every time the passive environment is exploited,
    from engineering high yield, low footprint farming to technologically harnessing the great and small Anomalies,
    including detoxification of materials and beings.
    \item \textbf{Blue} Biotechnology has developed from being used for aquatic Biotech, to include Megafauna, since the
    first Megafauna is told to having been spottet in the ocean.
    It allows scanning, categorization and assessment of Megafauna, including in some cases,
    harvesting Megafauna specialities.
    \item \textbf{White} or Industrial Biotechnology  is used in Industrial processes and in handling some of the
    results of such processes.
    Creating and applying contamination-sealant is one of such processes.
    \item \textbf{Black} deals with with integration of Biology and Technology from enhancements or replacements like
    Cyberware or more conventional Prosthetics to synthesizing life, ressurection and strong AI. Seldomly useful
    and the most dangerous brand of Biotechnology, it is not taught often.
\end{enumerate}
\subsection{Computer}\label{subsec:computer}
The Computer skill group describes the ability and training of the character to generally interact with computers.
\begin{enumerate}[label= -]
    \item \textbf{Programming} is creating or changing computer programs.
    \item \textbf{Usage} is using complex features of computers in intended ways.
    \item \textbf{Hacking} is making computers behave in unintended ways.
\end{enumerate}
\subsection{Engineering}\label{subsec:engineering}
The Engineering skill group describes the ability and training of the 
character to generally create and care for machines.
\begin{enumerate}[label= -]
    \item \textbf{Repair} 
    \item \textbf{Design} 
    \item \textbf{Build} 
\end{enumerate}
\subsection{Aux}\label{subsec:aux}
Auxiliary Training is training with specific types of equipment and tactics.
This includes operation and maintenance of these devices.
\subsubsection{Defense}\label{subsubsec:defense}
Skillgroup pertaining to defensive Equipment options.
\begin{enumerate}[label= -]
    \item \textbf{Decoys} From Flares against heat seeking missiles to Holographic Projectors, anything that guides
     an opponents Weapon to where it can inflict no harm.
    \item \textbf{Shields} Devices that dampen or stop kinetic energies, especially those of weapons
    \item \textbf{Deflectors} Anything that redirects energies, especially those of weapons
\end{enumerate}
\subsubsection{Offense}
Skillgroup pertaining offensive Equipment and Tactics
\begin{enumerate}[label= -]
    \item \textbf{Explosives} Training with things like mines, piles of C4 and IEDs
    \item \textbf{Sniping} Equipment and Tactics to specifically target enemy Systems and increase accuracy
    over rate of firing
    \item \textbf{Gadgets} Hidden weapons, fireable homing beacons, Flare guns, Flamethrowers, Cryothrowers, Gas
\end{enumerate}
\subsubsection{Recon}
Skillgroup pertaining to gathering information while remaining unseen
\begin{enumerate}[label= -]
    \item \textbf{Scanners} Equipment to obtain Information, from Sensors over Radar to Compositionscanners
    \item \textbf{Suppressors} Stealth equipment from mecha-weapon compatible flash or noise suppressors to thermal
    signature dissipating gadgets and even invisibility ``cloaks''
    \item \textbf{Sabotage} EMP, ECM, ECCM, but also poison, pitfalls and similar tactics
\end{enumerate}
\subsubsection{Utility}
Skillgroup pertaining to opening up more options
\begin{enumerate}[label= -]
    \item \textbf{Communications} Equipment with the purpose of transferring Information
    \item \textbf{Reprocessing} Equipment to turn things into other, generally more useful things.
    From electrolysis to cooking.
    \item \textbf{Utility Movement} Equipment that allows movement in nonstandard ways, such as flying, climbing,
    swimming, boring or even teleporting
\end{enumerate}
\subsection{Professions/Proficiencies}\label{subsec:professions}
each of these are a Skillgroup, meaning after the initial point, further progress needs to be specialized.
The exact nature of these specializations is open to player and storyteller.
\subsubsection{Metal}\label{subsubsec:metal}
\subsubsection{Ceramics}\label{subsubsec:ceramics}
\subsubsection{Leatherworking}\label{subsubsec:leatherworking}
\subsubsection{Boneworking}\label{subsubsec:boneworking}
\subsubsection{Woodworking}\label{subsubsec:woodworking}
\subsubsection{Woodcraft}\label{subsubsec:woodcraft}
\subsubsection{Animalcare}\label{subsubsec:animalcare}

\section{Spiritual Skills, Techside}\label{sec:spiritual-skills,-techside}
The Skills within Mecha Weapons Expertise and Piloting Practice are not grouped.
\subsection{Mecha Weapons Expertise}\label{subsec:mecha-weapons-expertise}
Mecha Weapons means any weapon too big to be carried, including stationary turrets.
Usually Weapons are fired with a combination between Weapon Type + Attack mode and \hyperref[itm:focus]{Focus}.
\subsubsection{Laser}
Laser or technically "Directed Energy Weapons" \par
are a group of weapons which typically do not use Ammunition, but raw Energy to inflict their damage.
Their comparably low damage is made up for by being as accurate as a light beam,
plus, the high amounts of particles in the air make for pretty lights.
\subsubsection{Projectile}
Projectile Weapons or sometimes just "Guns" \par
are a group of weapons which discharge Projectiles towards a target.
Ammunition has to be employed, but there are several different types available.
\subsubsection{Missile}
Missile Weapons are a weapons that travel to the target to inflict damage, the difference
to Projectile Weapons is that Missile Weapons are usually self propelled instead of being accelerated at their
startpoint.
Missiles are comparatively heavy and expensive, but most of them ignore shields and/or deliver
special devices or bombs to the target.
\subsubsection{Attack Modes}
\begin{enumerate}[label= -]
    \item \textbf{Direct} Weapons are many and varied but share the characteristic of being pointed roughly at
    the target they are shooting.
    They require line of sight.
    \item \textbf{Indirect} Weapons are just as varied as Direct Weapons but have the key difference
    of not pointing at the target.
    Usually Up and in an angle, but not always.
    They require information about the targets position, but are free from line of sight restrictions, as long as
    the Attackstill has a valid way to travel to the target.
    \item \textbf{Seeking} Weapons have some sort of sensors or a dataconnection, and will follow a Target, as long as the Lock
    remains unbroken.
    This skill mainly deals with the correct usage and parametrization of such Weapons.
    \item \textbf{Beam} Weapons that emit long, sometimes even sustained streams, making them easier to aim
    and inflicting more damage, usually at the price of higher, sustained costs and less flexibility.
\end{enumerate}
\subsubsection{Melee}
Melee Weapons of mech sizes are capable of inflicting utter destruction, but are usually
limited to close ranges, which makes them not very useful at long distances.
\begin{enumerate}[label= -]
    \item \textbf{Sharp} Melee Weapons have slightly longer range and are faster and usually do good damage.
    Most of them use the Movement System to get ready to attack again, so they may attack a lot of times on fast mechs.
    They also are the most precise of the melee weapons.
    \item \textbf{Blunt} Melee Weapons are usually slow, and easier to dodge, but their raw damage output is only
    matched by Artillery.
    \item \textbf{Unarmed} is not technically a category of weapon, but it represents being able to fight in melee
    without any special weapon, which allows for mostly disabling maneuvers.
\end{enumerate}
\subsection{Piloting Practice}\label{subsec:piloting-practice}
Piloting Rolls are usually done with a combination of Terrain+Weightclass and \hyperref[itm:discipline]{Discipline}

\subsubsection{Terrain}\label{subsec:terrain}
\begin{enumerate}[label= -]
    \item \textbf{Plain} Terrain
    \item \textbf{Desert} Terrain
    \item \textbf{Forest} Terrain
    \item \textbf{Arctic} Terrain
    \item \textbf{Mountain} Terrain
    \item \textbf{Urban} Terrain
    \item \textbf{Mud} Terrain
    \item \textbf{Marine} Terrain
    \item \textbf{Extreme} Terrain
\end{enumerate}
\subsubsection{Weightclass}\label{subsec:weightclass}
\begin{enumerate}[label= -]
    \item \textbf{Exo} is short for "Exoskeleton" and refers to Mechs that are barely larger than an extended Bodyarmor
    or Mechs that are generally not much bigger than their Pilot. (<1t)
    \item \textbf{Tiny} a versatile spot for Contraptions between Exo and Very Light, these include buggies, cars and
    a lot of small, fast Vehicles. (1-5t)
    \item \textbf{Very Light} include things like APCs and small fast walkers and light tanks.(5-15t)
    \item \textbf{Light} is about as big as a medium to main tank. (15-50 t)
    \item \textbf{Medium} is about as big as a heavy to superheavy tanks (50-150 t)
    \item \textbf{Heavy} is above superheavy tanks. (150-500t)
    \item \textbf{Very Heavy} LANDKREUZER (500t-1500t)
    \item \textbf{Ultra} ???
\end{enumerate}


\section{Spiritual Skills, Mageside}\label{sec:spiritual-skills,-magicside}
Magic is to be determined with the magicuser.

\subsection{Raven}\label{subsec:raven}
This skillgroup is an example of a Tribal Raven shamans link with their totem
\begin{enumerate}[label= - ]
    \item \textbf{Talon} generally damaging and interacting spells
    \item \textbf{Eye} scouting and information spells
    \item \textbf{Wing} movement and utility spells
\end{enumerate}

\subsection{Spirit}\label{subsec:spirit}
This skillgroup is an example of a Tribal Raven shamans general ability to manipulate Contamination
\begin{enumerate}[label= - ]
    \item \textbf{Call} increases Contamination at target (further away => weaker effect)
    \item \textbf{Infuse} increases Contamination in target and restores it
    \item \textbf{Expunge} decreases Contamination in target and damages it
\end{enumerate}


\subsection{Resonant Body}\label{subsec:resonant-body}
This skillgroup is an example of a Tribal Warriors awoken innate superhuman abilities.
\begin{enumerate}[label= - ]
    \item \textbf{Resonant Healing} Every 10-Discipline (maximum Resonant Healing) Contamination
    provides a bonus die when healing.
    \item \textbf{Tireless} Focus negates exhaustion, requires Contamination lasts 1/4/16 hours per day
    \item \textbf{Ancestors Strength} Intuition Adds to Physical skill in times of need (Ancestors Strength times per day)
\end{enumerate}

\subsection{Deathless}\label{subsec:deathless}
This perkline provides the legendary resilience seen in some Tribals
\begin{enumerate}
    \item \textbf{Last Chance} No Single Wound can knock the character out
    \item \textbf{Return} If the Body is restored soon after death, it will come back to life.
    \item \textbf{Join the Ancestors} On Death Character becomes a ghost and is able to communicate with receptive minds.
\end{enumerate}


\subsection{Generalized, Vague, Pointsink}\label{subsec:GVPS}




Attributes represent in what aspects of the gameworld a character excels and in which he doesnt.\par
Depending on the Priorities picked during character creation it may be easier or harder to raise them, but every
Attribute can be Applied to a multitude of general Actions, see the examples in the table below.\par
The Priorities and Attributes in the following Table are a not all that exist
but they are the most common Attributes shared by all humans.\par
\begin{tabular}{l||cccccc}
    Priorities      &Social        &Mental         &Physical    &Ability     &Spiritual  \\\hline\\
                    &Resolve       &Perception     &Fitness     &Competence  &Intuition \\
    Attributes      &Impression    &Mind           &Dexterity   &Insight     &Focus\\
                    &Sympathy      &Knowledge      &Agility     &Theory      &Discipline\\



\end{tabular}\vspace{1cm}
\section{Physical Attributes}\label{sec:physical-attributes}
\begin{enumerate}[label= -]
    \item {Fitness}: general physical actions, applying strength, persisting, healing
    \item {Dexterity}: precise and/or speedy manipulation of objects using hands
    \item {Agility}: running, jumping, dodging, climbing
\end{enumerate}
\section{Social Attributes}\label{sec:charisma-attributes}
\begin{enumerate}[label= -]
    \item {Resolve}: reasoning, convincing, bartering, negotiating and similar interactions
    \item {Impression}: leadership, bluffing, acting, boasting, lieing
    \item {Sympathy}: chatting, pleading, resolving conflicts, being likeable
\end{enumerate}
\section{Mental Attributes}\label{sec:mental-attributes}
\begin{enumerate}[label=-]
    \item {Perception}: gathering information, noticing details, being alert
    \item {Mind}: memory, assessing situations, processing information, menial mental tasks
    \item {Knowledge}: knowing, learning, possessing and storing information
\end{enumerate}
\section{Ability Attributes}\label{sec:ability-attributes}
\begin{enumerate}[label= -]
    \item {Competence}: repetetive or physically complicated actions
    \item {Insight}: creative or explorative actions
    \item {Theory}: applying knowledge or mentally complicated actions
\end{enumerate}
\section{Spiritual Attributes}\label{sec:spiritual-attributes}
\begin{enumerate}[label= -]
    \item {Intuition}\label{itm:intuition}]: knowing without thought
    \item {Focus} \label{itm:focus}: the power of concentrated effort
    \item {Discipline} \label{itm:discipline}: to do things many times done before
\end{enumerate}
\section{Character Attributes}\label{sec:character-attributes}
\begin{enumerate}[label= -]
    \item {Contamination Resistance}: see \hyperref[subsec:categories]{Contamination Categories}
    \item {Technological Affinity}: dictates the Techlevel that can be used with minimal risk
    \item {Luck}: the general tendency towards good or bad random outcomes
\end{enumerate}


\chapter{Trade}\label{ch:trade}
These entries in the table below are usually equivalent, local supply and demand may vary them greatly.
Rolling for price variance might be done if there is no in-world reason why certain things might be more or
less expensive.
If that is the case the prices vary by \(((-1)^{1d10}\cdot1d10\cdot5) \%\).
Additionally everyone a character trades with has self interest and will keep some of
the tradevalue.\par
\begin{tabular}{cccccccc}
threshhold&0&5&8&11&14&17&20\\
kept tradevalue&50&33&20&10&5&2&0
\end{tabular}\\
Goods, at 100\% value each line is equivalent to one another and equivalent to 4 character creation points \par
\section{Tradegoods}\label{sec:tradegoods}
\begin{tabular}{|r|l|}
    \hline
    Storage & Article\\\hline
    10g & Elixir of Life (EOL)\\\hline
    100g & Alacast\\\hline
    500g & Experimental Tech Scrap (ETS)\\\hline
    1kg & High Tech Scrap (HTS)\\\hline
    1.5kg & Advanced Gene Therapy Medicine (AGT)\\\hline
    5kg & Medium Tech Scrap (MTS)\\\hline
    10kg & Potent Medicine\\\hline
    50kg & Low Tech Scrap (LTS)\\\hline
    50kg & Seeking Rockets\\\hline
    100kg & Processed Medicine\\\hline
    100kg & Basic Rockets\\\hline
    100kg & Artillery\\\hline
    200kg & Basic Medicine\\\hline
    500kg & Improvised/Herbal Medicine\\\hline
    500kg & Base Tech Scrap (BTS)\\\hline
    500kg & Basic Ammunition\\\hline
    1000L & drinking Water (500days of drinking)\\\hline
    500L/35kg & Liquid Hydrogen\\\hline
    1t & Coal\\\hline
    350L & LiquidCombustionFuel\\\hline
    100kg & High Energy Rations (0.4kg/(day\(\cdot\)person))\\\hline
    1t & Normal Food (2.5 kg / (day\(\cdot\)person))\\\hline
    1 & Medium Quality Blueprint\\\hline
\end{tabular}\par
Open for more Suggestions\vspace{1.5cm}
Example:\par
Character A wants to get rid of 5 tonnes of Base Tech Scrap in favor of more easily transportable Alacast in a local
Dome.\par
If there are no storyelements influencing the prices, they are rolled.
First, the value of the BTS is determined to be (Roll: 6, 4) \(((-1)^{6}\cdot4\cdot5) \% = +20\%\).
Then the value of Alacast is determined
to be (Roll: 10, 3)\(((-1)^{10}\cdot3\cdot5) \% = +15\%\).
The result is that the rate between BTS and Alacast is
\( 500\cdot1.15/100\cdot1.2 = 4.79\) \\ Which means character As scrap is worth
\(5000 kg\cdot / 4.79 = 1043 g \) of Alacast.
He barters (with Resolve and Trade for 3, 2) with a local merchant and manages to
negotiate terms in which the merchant retains (Roll: 2, 4, 5, 8, 9 \(\Rightarrow 9 \Rightarrow 10\%\) of the Traded
value.
He decides to trade and gets 938 grams of Alacast.\vspace{1.5cm}
\pagebreak



\chapter{GameSystems}\label{ch:gamesystems}
\section{Health and Contamination}\label{sec:healthandcontamination}
\subsection{Healing}\label{subsec:healing}
For living entities damage is represented by wounds with the severity being the remainder of the damage after all
defenses were subtracted.
Humans roll a Fitness-Check once a wound starts healing.
For every met threshhold of the wound one level of regeneration rate for that particular wound is noted down.
If a wound is worsened/increased the Fitness-Check is repeated, selecting the worse of the two results.
If a wound is successfully and substantially improved, the Fitness Check is repeated, selecting the better of the two
results.
Regeneration can be aided or slowed by circumstances, such as therapy or environmental conditions.
Resonance of frequency 1 lowers the healing rate by its amplitude, even going negative, while resonance of frequency 10
improves healing rate by its amplitude.
Wounds do not necessarily lead to death or permanent impairments, but if they remain untreated for a long time, or there
is a runaway effect on one of them, the Storyteller may decree permanent repercussions, as appropriate, including Death.
Regeneration accumulates over successive days, and once the current severity is reached, the wound is lowered by 1
severity and regeneration is reset to 0.
If regeneration rate is negative, and the regeneration rate falls below 0, regeneration progress is set to the wound
severity and then the severity is increased by 1.
Standard threshholds are 2, 4, 6, 8, 10, 12, 14. \par
The following table lists healing time in days for severity and hit threshholds.\par
\begin{tabular}{c|rrrrrrrrrr}
    Severity & 1 & 2 & 3 & 4 & 5 & 6 & 7 & 8 & 9 & 10\\
    1 &  1 d & 1 d & 1 d & 1 d & 1 d & 1 d & 1 d & 1 d & 1 d & 1 d\\
    2 &  3 d & 2 d & 2 d & 2 d & 2 d & 2 d & 2 d & 2 d & 2 d & 2 d\\
    3 &  6 d & 4 d & 3 d & 3 d & 3 d & 3 d & 3 d & 3 d & 3 d & 3 d\\
    4 &  10 d & 6 d & 5 d & 4 d & 4 d & 4 d & 4 d & 4 d & 4 d & 4 d\\
    5 &  15 d & 9 d & 7 d & 6 d & 5 d & 5 d & 5 d & 5 d & 5 d & 5 d\\
    6 &  21 d & 12 d & 9 d & 8 d & 7 d & 6 d & 6 d & 6 d & 6 d & 6 d\\
    7 &  28 d & 16 d & 12 d & 10 d & 9 d & 8 d & 7 d & 7 d & 7 d & 7 d\\
    8 &  36 d & 20 d & 15 d & 12 d & 11 d & 10 d & 9 d & 8 d & 8 d & 8 d\\
    9 &  45 d & 25 d & 18 d & 15 d & 13 d & 12 d & 11 d & 10 d & 9 d & 9 d\\
    10 &  55 d & 30 d & 22 d & 18 d & 15 d & 14 d & 13 d & 12 d & 11 d & 10 d\\
    11 &  66 d & 36 d & 26 d & 21 d & 18 d & 16 d & 15 d & 14 d & 13 d & 12 d\\
    12 &  78 d & 42 d & 30 d & 24 d & 21 d & 18 d & 17 d & 16 d & 15 d & 14 d\\
    13 &  91 d & 49 d & 35 d & 28 d & 24 d & 21 d & 19 d & 18 d & 17 d & 16 d\\
    14 &  105 d & 56 d & 40 d & 32 d & 27 d & 24 d & 21 d & 20 d & 19 d & 18 d\\
    15 &  120 d & 64 d & 45 d & 36 d & 30 d & 27 d & 24 d & 22 d & 21 d & 20 d\\
    16 &  136 d & 72 d & 51 d & 40 d & 34 d & 30 d & 27 d & 24 d & 23 d & 22 d\\
    17 &  153 d & 81 d & 57 d & 45 d & 38 d & 33 d & 30 d & 27 d & 25 d & 24 d\\
    18 &  171 d & 90 d & 63 d & 50 d & 42 d & 36 d & 33 d & 30 d & 27 d & 26 d\\
    19 &  190 d & 100 d & 70 d & 55 d & 46 d & 40 d & 36 d & 33 d & 30 d & 28 d\\
    20 &  210 d & 110 d & 77 d & 60 d & 50 d & 44 d & 39 d & 36 d & 33 d & 30 d\\
    21 &  231 d & 121 d & 84 d & 66 d & 55 d & 48 d & 42 d & 39 d & 36 d & 33 d\\
    22 &  253 d & 132 d & 92 d & 72 d & 60 d & 52 d & 46 d & 42 d & 39 d & 36 d\\
    23 &  276 d & 144 d & 100 d & 78 d & 65 d & 56 d & 50 d & 45 d & 42 d & 39 d\\
    24 &  300 d & 156 d & 108 d & 84 d & 70 d & 60 d & 54 d & 48 d & 45 d & 42 d\\
    25 &  325 d & 169 d & 117 d & 91 d & 75 d & 65 d & 58 d & 52 d & 48 d & 45 d\\
\end{tabular}
\subsection{Treatment}\label{subsec:treatment}
Treating of wounds is usually done with a Competence or Theory + Healing Check, where medicine and equipment
can provide modifiers or advantages.
Many wounds require medicine to be treated at all.
Treating them without will either not be possible or generate disadvantage.
Standard threshholds are 8, 11, 14, 17, 20, with results below 5 having negative consequences.

Treating ingress wounds requires at least 10g of Alacast, with the treatment of
contamination inflicting the sum of all dice including and below or equal to a threshhold determined by the
Technology used (i.e.\ 8 for Saline-Alacast-Solution Injection) as damage.
The level of ingress wounds and contamination is directly lowered by the number of hit threshholds.

In all cases Alacast is used to precipitate Contamination from the body, the
crudest way is to introduce Alacast to the bloodstream to precipitate it inside the body and hope the body expels
it on its own, more refined ways are basically a dialysis or a very specific targetting of crystallization seeds
and retrieval of precipitation clusters.

\subsection{Medicine}\label{subsec:medicine}
Medicine is used while making a Treatment Check.
Each Treatment of a wound using Medicine consumes
severity\(\cdot\)100g of the medicine and provides the specific bonus.

\begin{enumerate}[label= -]
    \item \textbf{Naturopathy} uses healing and the respective field of the medicine.
    It provides a good floor and good threshhold bonuses.
    \item \textbf{Specialized Medicine}(name pending) live mixing of specialized cures from ingredients
    uses healing and the respective field of the ailment.
    It has no floor, but scales the best.
    \item \textbf{Broadband Medicine} is mass produced, refined medicine and uses a com\-bi\-nation of
    Healing, \-Scie\-nce(Med\-icine) and \-Red Biotech.
    It has a good \-floor, but low \-scaling.
    \item \textbf{Alacast} has no direct medicinal benefit, but effectively combats Contamination in all fields
    \item \textbf{Elixir of Life} is the top notch medicine, rejuvinates and heals nearly anything.
    Sadly, it contaminates the user.
\end{enumerate}

\subsection{Categories}\label{subsec:categories}
Every Character has a Contamination Resistance rating ranging from 1 to 5,
but in-world they are categorized (by supervivo) from A to C\@.
\begin{enumerate}[label = - ]
\item \textbf{CatA} Humans are usually Supervivo and usually have low Contamination Resistance
                    and high Affinity.
\item \textbf{CatB}s are outlaws, traders, hermits or for some other reason living outside the Domes and Tribes.
                    They can have any, but usually have medium Contamination Resistance and medium Affinity
\item \textbf{CatCs} are usually Tribals with high Contaminatin Resistance and low Affinity.
\end{enumerate}
\begin{tabular}{c|cc}
    Entity & Contamination Resistance  & Affinity requirements\\
    Experimental & 1 & 5 \\
    High Tech & 2 & 4\\
    Mid Tech & 3 & 3 \\
    Low Tech & 4 & 2\\
    Base Tech & 5 & 1\\
\end{tabular}

\textbf{Affinity requirements} are the base amount of Affinity that a character has to have to use the Entity.
If the requirements are unmet, every use incurs the difference in penalty dice.


\subsection{Contamination}\label{subsec:contamination}
Contamination is a term used to describe the amalgamation of \-technology destroying \-nanoweapons, \-radiation, \-pollution
and general \-environmental \-hazards.
Contamination is air\-borne, waterborne, bloodborne and per\-meates pretty much everything in the environment.
It is usually assigned a \-level \-describing its intensity.
Only within the Domes of the Supervivo, \-inside a Mech or in a similarily \-purified environment can the \-
Contamination\-level ever reach 0. \par
Levels below 0 are necessary for taking apart Experimental Tech Technology, but getting an area so pure is hard.\par
\begin{enumerate}[label = - ]
\item Levels 1--3
are rare areas of low Contamination, like mines, airlocks, outdoor markets, outdoor settlements and so on.
Even subjected to this level of Contamination for months, it will not cause death.
However, for the more fragile Members of Society, the general quality of life can be greatly diminished.
\item Levels 4--6
are the most common and usually inhabited by Cat-B\@.
People without a form of permanent resistance can die here, but life expectancies should be a few months at least.
\item Levels 7--11
are usually inhabited by tribals but sometimes a few Cat-B have to make Camp here.
Anyone without contamination protection will die here within a few weeks.
\item Levels 12--15
are sparsely populated by hardy Tribals, but usually deserted.
\item Levels 16+
are the most hostile areas.
Weird Things happen here, but some say that there is great treasure\ldots
\end{enumerate}
There is no upper limit for contamination levels.
If an entity is in an area with contamination, every hour it receives the local contamination
as contamination damage.\par

\subsection{Contamination damage}\label{subsec:contamination-damage}
Anytime the an entity receives contamination damage, the character lowers that damage with a roll on their contamination
resistance and fitness plus all applicable modifiers (external only for contamination ingress).
Remaining Damage above 0 is applied to the ingress wound (extending the existing ingress wound if possible,
otherwise creating a new one) and raises contamination of that character by 1.
The contamination of the character is applied as an internal bonus on the contamination resistance result, but reduces
their healing rate, including going negative.



\subsubsection{Contamination Damage to Entities}
When a character directly interacts with an entity, the Charactercontamination\(\times\)2
of that character is applied as Contamination damage.
This is repeated every hour if the interaction continues.
While directly contacting an entity only the internal resistance applies, direct contact might be piloting a mech,
using a gun or shaking someones hand, except when proper seals are in place.

Instead of Fitness an object uses how well adapted it is to being in contact with contamination.
This is a combination of relevant factors.
In general the more artificial an Object is the easier it is broken down,
the more regular it is the easier it is disturbed and the more complicated it is the easier it is to bring into dissaray.

Some Examples would be

\begin{enumerate}
    \item Rock/Stick/Clump of Dirt: might bear corruption but is usually not affected
    \item Bronze Sword (5,5): the Edge will dull in days or weeks and the material will acquire a thick Patina,
    but over all it is quite usable and with the patina sealing it, it would last a long time.
    \item Combustion Motor(4,4): a few hours in Corrupted air and it will deteriorate
    \item Vaccuum Tubes(3,4): from the perspective of Corruption ``just'' glass, but all the interesting parts around it
    \item High Tech Supervivo Survival suit(2,5): Made from simple, chaotic polymers, it can
    withstand Corruption even without a Sealant for several days before crumbling
    \item 3nm Hightech Processor(0,2): Vulnerable on every front.
\end{enumerate}


Damaged Objects might loose some or all of their functionality, or cause penalty dice when using them.

\subsection{Seals}\label{subsec:seals}
A seal of a level of X lowers the level of the Contamination behind it by X
If a seal is damaged, its effectiveness is lowered by 1 and until it is patched, it only counts for half
its current value

Air has to be purified separately or enriched with oxygen.
(One human consumes about 500--600 litres of oxygen per day.)
The Domes of the Supervivo for
example usually have at least a level 25 hermetical seal with purified and oxygen enriched air.

Examples of human sized external contamination modifiers or in-place options:

\begin{tabular}{cl}
    1& An alacast infused rag or a shut wooden door (although oxygen may run out)\\
    2& a breathermask or a simple wooden door airlock (again, oxygen)\\
    3& a partial gasmask or an airtight airlock\\
    4& a heavy partial gasmask with Goggles or a basetech air filtering facility\\
    5& a full gasmask\\
    6& as above + sealed clothing or lowtech air filtering facility\\
    7& as above + midtech scrubber (a little backpack)\\
    8& as above + hightech air scrubber or midtech air filtering facility\\
    9& as above + full mask with air from a compressed air tank\\
    10+& \hyperref[subsec:seals]{specially sealed} suit with airtank or hightech filtering facility\\
\end{tabular}


\section{Computers and Programming}\label{sec:computers}
\subsection{Computers}\label{subsec:computers}
Computers come in many different shapes and sizes, but in Endworld they are all abstracted into
their availaible memory and their available processing powers.

Computers in Endworld are different from our own, since they had to be designed to be much easier, much more
modular and much more resistant, they are usually quite big and heavy, but also incredibly sophisticatedly
parallelized.
Also they are heavily abstracted and should not be confused with real computers.

\subsubsection{Processing Power}
Processing power is a computers capability to achieve computation and the general work, including I/O\@, it is measured
in Computational Power Usage, and a Processor provied CPU capacity.
Processing power is occupied by programs, taking as much as they need or failing to start if there is not enough.
A Theory, Computer Usage (and relevant skill for computer administration, clocking or similar) Roll can be used to shift 
[8,12,17,20] processing from one program to another or free that processing power up.
A program running below its required processing power, will perform worse.
Scale any applicable metric to the utilized processing power, rounding unfavorably for purposes of the program.
Programs with no direct metric, still produce lesser effects, or may not work at all.
A program running above its required processing power usually does not do its tasks better, but might perform them
faster.

\subsubsection{Storage Space}
Storage space is a computers capability to store data and programs.
It is measured in Storage.

\subsubsection{Programs}
For any Program to run, it needs to either be loaded from storage or be present as built in RROM\@.

\paragraph{RROM}
Robust Read Only Memory is a storage medium directly integrated into a processor and by custom usually made from non
volatile, corruption resistant materials like woven strands of wire and beads.
Since a program is needed to load a program, some processors have built in RROM acting as a bootstrap loader for the
Main program used to load other programs more efficiently.
Since the RROM is Using its own Computational Power it can not be reused and the Bootloader is effectively its own
Computersystem.
It starts executing as soon as it is turned on and the program can not be changed or modified while the Processor is
running.
However simply physically replacing the program is enough and so RROM has a big advantage over hardware or even elektro-
mechanic computing techniques.

\paragraph{Execution}
Programs have a set of goals that they will try to accomplish and through their programming take steps towards.
This culminates in the abstracted behaviour of programs performing actions, some with a predetermined degree of success,
some provide parts or all of the statistics needed for a check.
Programs can load other programs either in parallel (while they are still running) or in series
(as they are completing).
When a program is completed, it unloads and at the end of the turn any ressources it used are freed.

\paragraph{Operating System}
An Operating system or ``Main'' is simply another program that can be dynamically tasked by the user to load other
programs.
A user interaction requires usually one turn, after which the Operating system performs a predetermined series of loads
and supplies each program with its parameters.

\paragraph{Parameters}
Are the context a specific program needs to accomplish its goal.
Parameters are usually abstracted away as long as it is reasonable to assume that all programs ``know'' what they
are acting upon.

\subsubsection{Programming}
To create a Program requires some knowledge of what the Program is trying to accomplish, and Computer Programming and
usually Theory.
If the process is more exploratory and there is no clear goal, Insight is used instead to merely
determine if a solution can be found.
For Theory checks, given that the programmer could conceivably find a series of Instructions, The Storyteller and
Player determine a set of CPU requirements, Storage space needed for the Program, auxiliary data requirements, flat
Checkresults or purposeful statistics, other programs loaded and runtime, as far as applicable for that program, with
similar Programs and their checked quality as templates.
After the Check was made, the Storyteller offers several configurations of the above values that could be achieved with
the roll and the Player chooses one.
Alternatively the Player designs a new Program, its configuration and effects and the Storyteller tells the Player the
interval of the programming rolls and the total sum of checkresults needed to achieve the desired goal.
The program, its name, its statistics and effects are then noted down by the Player.

\paragraph{Example - Shield Daemon}
\begin{enumerate}
    \item Name: Shield Daemon
    \item CPU: 0.5
    \item Size: 1
    \item Requirement: Internal Computerization, Shields
    \item Effect: If a switch inside the cockpit is flipped to ``bring shields up'',
    load one Shield Configurator for each shield that is down,
    if it is flipped to ``NOW'', load one Shield Burster for each shield that is down
    \item Runtime: Ongoing
    \item To program: Shields or relevant piloting skill, 1 workday per roll, total: 100
    \item Limits: specific to exact programs and mech
\end{enumerate}

\paragraph{Example - Shield Configurator}
\begin{enumerate}
    \item Name: Shield Configurator
    \item CPU: 3
    \item Size: 1
    \item Requirement: Shield
    \item Effect: Configure Shield with Flat Result of 10
    \item Runtime: 1 Turn
    \item To program: Shields, 1 workday per roll, Any 3 Rolls of 16+
    \item Limits: specific to exact shield
\end{enumerate}

\paragraph{Example - Shield Burst}
\begin{enumerate}
    \item Name: Shield Burster
    \item CPU: 1
    \item Size: 1
    \item Requirement: Shield
    \item Effect: Configure Shield with Flat Result of 0
    \item Runtime: Instantly
    \item To program: Shields, 1 workday per roll, Total of 30+
    \item Limits: specific to exact shield
\end{enumerate}








\chapter{Combat}\label{ch:combat}
\section{Rounds}\label{sec:combat-rounds}
Combat is done in rounds.
In Detailed Combat, each round is 5 seconds long, in Tactical Combat, rounds are of variable length and could
last an hour.
Every character present acts in turn (determined by circumstances or a Stealth, Perception, Willpower or other
situational check) and can usually move and act or attack.
Double actions, if possible, usually carry a 2 dice penalty on both actions, further actions carry even more penalties.

\section{Attacks}\label{sec:combat-attacks}
An attack has to be made on a valid target.
Targets are valid depending on what weapon is used, a direct fire weapon requires line of sight, for example.
Depending on the circumstances, the correct attack roll is made, usually involving dexterity, the weapon skill and
another skill depending on circumstance.
This roll, called attack value from here, is then modified by range depending on the weapon.

If the attack value is positive, the general vicinity of the target is hit.
Then, the enemy defense from dodging or cover is subtracted.

If the attack value is still positive, the target is hit and Damage is potentially inflicted,
Otherwise, the cover may still be damaged, depending on circumstance.


\section{Damage}\label{sec:combat-damage}
Entities have several hitzones, which zone specifically takes the is determined by the Entity and potential targeting.
Every successful hit will damage the hitzone, eventually leading to destruction.
A hits attack value is mitigated by armor.

If the remaining strength is above 0, the hitzone is damaged.

If it is under 0, the hitzone may still count as scratched.

Layered defenses do not provide protection if they are destroyed, some may already "leak" damage onto the sublayer
if they get damaged.

For a Biological Entity damage represents wounds and may heal.


\section{Death}\label{sec:combat-death}
There is no fixed numerical amount of life points, but every time a hitzone, including body parts, is damaged,
this damage is represented by some form of reduction in operation.
This can manifest in many ways, like penalty dice for associated activities, higher botch/malfunction frequencies,
lowered armor, and many more.
If there is a lot of damage, machines may become inoperable, people might pass out.
Death and Destruction only occur when the game master decrees so.

\section{Defense}\label{sec:combat-defense}
Defense can be gained from dodging or cover or other situationally appropriate means.
In general the defense value is just the roll of the action, if it would provide appropriate protection.
Taking cover behind a hanging fishing net or trying to talk down an automated turret might not result in any defense
at all.

\textbf{Dodging}, usually done with agility and acrobatics, footwork or running, requires the action for that round,
meaning it needs to be repeated every round, and is subject to double action rules.

\textbf{Cover}, once taken with agility and tactics, running, stealth, footwork or whatever is appropriate,
provides its defense until it is nullified by flanking or movement is taken.
Furthermore, there is the possibility of full cover granting a flat bonus onto the roll

\textbf{Other Actions} may include throwing sand, building rapport, talking someone out of shooting and are resolved
by themselves


\section{Detailed Combat}\label{sec:detailed-combat}
\begin{enumerate}
    \item Every participant (technically in secret) decides on what they are going to do.
    \item Non-offensive actions are resolved first, then offensive actions in descending order of rolled result.
\end{enumerate}



\section{Tactical Combat}\label{sec:tactical-combat}
In many respects much like detailed combat, Tactical Combat zooms out over encounters that could take a while,
it is not unusual to switch back and forth.
In Tactical Combat, there are usually no single actions, Tactical Actions encompass states and processes.
Anytime something changes, everyone who is aware may change what they are doing.
If these reactions cascade, detailed combat is the natural consequence and actions will proceed in 5 second intervals.


If an actor does more or less the same thing 3 times in a row, it might be appropriate to enter Tactical Combat.
A state and/or process will interpret the last 3 rolls and take the average of that number.
Rolls that do not succeed fully may carry over some part of the result until the threshold is met, so someone nearly
missing will instead hit at a lower rate


When entering Tactical Combat some form of time commitment is agreed upon and after each interval one roll is made.
The moment of this roll is also the only time for unprovoked re-evaluation and change of plans.
The interval should be changed by the Storyteller depending on the density of action.
Tactical Actions might include \par
\begin{enumerate}
    \item firing on an enemy or a position
    \item guarding a position
    \item breaking down a door
    \item traversing to a point
    \item looking out
\end{enumerate}

\section{Targeting}\label{sec:combat-targeting}
Generally a center-mass targeting is assumed.
From there deviations occur, spreading the damage over neighboring parts
For a Human target, the hitzones are determined by d10, or by the resonance of the attack roll if any.
If there are several resonances, the attacker gets to pick.


\begin{tabular}{cc}
    roll & target\\\hline
    1,2,3 & torso \\
    4,5,6 & legs \\
    7,8 & arms \\
    9 & hands or feet \\
    10 & head
\end{tabular}

\subsection{Called Shots}\label{subsec:combat-calledshots}
To ensure that a specific area is hit, a called shot can be made, risking missing altogether.
A Penalty to the attack value is taken, to hit the specific body part.
A resonance with a number corresponding to a body part less than what was aimed for,
will instead attack that part, as above.

\begin{tabular}{cc}
    target & penalty\\\hline
    head & 3 \\
    hands or feet & 2 \\
    legs & 1 \\
    torso& 1
\end{tabular}

\section{Weapons}\label{sec:weapons}

Weapons have the following stats:
\begin{enumerate}
    \item \textbf{Range} - Minimum / Medium / Drop, in meters
        \begin{enumerate}
            \item \textbf{Minimum} - below which the weapon is hard to use in combat and penalty dice are issued
            \item \textbf{Optimal} - below which no penalties for range are in effect
            \item \textbf{Drop} - every how many meters a -1 penalty on shots is incurred
        \end{enumerate}
    \item \textbf{Ammunition} - how many shots per reload/how much weight per reload
    \item \textbf{Shot Interval} - in seconds,  values below 5 allow multiple shots per round
    \item \textbf{Handling} - initiative bonus
    \item \textbf{Mods} - as per mod
    \item[Crit] - requisites and effects of a critical hit
    \item \textbf{Skill} - the main skill to use this weapon
\end{enumerate}

\subsection{Example weapons used in Play}\label{subsec:example-weapons-used-in-play}
To be categorized and refined later.
These Weapons are of Personal Size.
\subsubsection{Sniper Rifle}
200[7,9,11]10\par
500[5,7,9,11,13,15,17,19]5\par
1500[7,9,11,14,17,20]2\par
6000[15,20]0.5\par
Costs: 1.5TU\par
Damage: 20\par
\hyperref[subsec:hit]{\textbf{Hit Calculation}} \par
[0,80,100,150]
Ammo 100g/shot

\subsubsection{Heavy Revolver}
3[5,7,9,11,13]2\par
20[5,7,9,11,13,15,17]10\par
50[9,14,18]5\par
100[12,16]5\par
150[18]3\par
Costs: 0.1TU\par
Damage: 15\par
Ammo 50g/shot

\subsubsection{Light Revolver}
3[5,7,9,11,13]2\par
20[5,7,9,11,13]10\par
50[7,12,16]5\par
100[14,17]5\par
Costs: 0.1TU\par
Damage: 10\par
Ammo 20g/shot

\subsubsection{Shotgun Revolver}
Slugs:\par
3[3,5,7,9,11,13]10\par
15[5,7,9,11,13]10\par
25[7,12,16]5\par
50[14,17]5\par
Shot:\par
3[3,5,7,9,11,13]10\par
36[5,8,11,14,17]
\hyperref[subsec:hit]{\textbf{Hit Calculation}} \par
[10,20,30,40,50,60,70,80,90,100]
Costs: 0.4TU\par
Damage: 20\par
Ammo 20g/shot


\subsubsection{Wooden Bow}
20[5,7,9,11,13,15,17,19,20]3\par
50[7,10,13,16,19]2\par
100[10,15,20]1\par
200[15,18,20]0.1\par
Costs: 0.1TU\par
Damage: 10\par
\hyperref[subsec:hit]{\textbf{Hit Calculation}} \par
[0,80,100,150]
Ammo 100g/shot

\subsubsection{Tooth Spear}
3[5,7,9,11,13,15,17,19,20]\par
Costs: 0.2TU\par
Damage: 30\par

\subsubsection{Obsidian Knife}
1[5,9,13,14,15,16,17,18,19,20]\par
Costs: 0.2TU\par
Damage: 20 \par
\hyperref[subsec:hit]{\textbf{Hit Calculation}} \par
[0,1,10,20,80,100,150]






\chapter{Tech}\label{ch:tech}
Electrolysis splits 1 L H2O into 622.22 L Oxygen und 1244.4 Hydrogen using 3.7037 kwh\\
\section{Engineering}\label{sec:engineering}
\subsection{Blueprints}\label{subsec:blueprints}
Blueprints are required for Engineer(Build) tasks.
To create a blueprint a character needs the appropriate materials, usually a computer of some sort, but even sand and a
stick can work.
The time required is based on the complexity of the project.
To find out if the blueprint is within the capabilities of the character check their specific knowledge + Theory
against the difficulty of the project.
If at least one threshhold is hit, the blueprint is viable.
Starting quality of a blueprint can be negative.
Every threshhold hit increases the quality of the blueprint by 1.
If the character does not have the Appropriate Knowledge he can check sufficiently similar knowledge with
an appropriate penalty dice (1 per techlevel, 3 for different module classes on the same techlevel)
Even a blueprint with a negative bonus is useful, since it enables building the item in question at all.\\
A blueprint states the name of the Item, its category and the required materials.

\subsection{Building}\label{subsec:building}
Building something requires a blueprint, the raw materials and the necessary tools.
Check Engineering modified by the blueprint and the conditions.
The finished product is [check result \(\cdot\) 10\%] Efficient.



\chapter{Mechs}\label{ch:mechs}

\subsection{Sizes}\label{subsec:sizes}

\begin{tabular}{r|ll}
    Sectors & Mechclass & Weight\\
    0.5 &Human & 0.1\\
    1 &Exo& 0.6\\
    2 &Exo,Tiny& 1\\
    3 &Tiny & 2\\
    4 &Tiny, Very Light & 4\\
    5 &Very Light& 10\\
    6 &Very Light, Light& 18\\
    7 &Light& 40\\
    8 &Light, Medium& 70\\
    9 &Medium& 150\\
    10 &Medium, Heavy& 300\\
    11 &Heavy& 500\\
    12 &Heavy, Very Heavy& 700\\
    13 &Very Heavy& 900\\
    14 &Very Heavy& 1200\\
    15 &Very Heavy, Ultra& 1500\\
    16+ &Ultra& >1500t\\
\end{tabular}\par
Modules are assigned to sectors.
Each sector has a size of 1, a module takes up some of that size.
Modules can be spread over multiple Sectors, as long as they are adjacent, although this will increase the likelyhood
of it sustaining damage.
Some Modules require a Hardpoint, there is generally only one Hardpoint per Sector.\par
A Mech has an effective Tech Level of the Tech at the 90\% point of the sizes modules ordered by Tech Level.\par
Contamination Damage applies to the highest Tech modules first.
Each Module can be damaged individually.\par

\section{Systems}\label{sec:systems}
A Mech needs a movement system, energy and heat systems and should probably have a seal, weapons and support systems.
\subsection{Costs}\label{subsec:costs}
Systems Examples\par
\begin{tabular}{c|c|l}
    Tech    & Credits Cost    & Examples\\
    B    &    4        &    \makecell*[{{p{12cm}}}]{Water Destillery, Rail based movement system, simple wheels,
    Manned gun, Ramshield, Floatation,   Metal Armorplating , Extra Cargo (10), steam engine, passive cooling}\\&&\\
    L    &    5        &    \makecell*[{{p{12cm}}}]{Lowtech Detox-Decon Equip(LDDE; +1), suspended wheels, tracks,
    bipedal, Dynamo , Autoturret, Jumpjets, Mech sized Sword, water cooling circle, basicHeatsink }\\&&\\
    M    &    6        &    \makecell*[{{p{12cm}}}]{Shields, flexible/balanced bipedal, Boostjets,
    Deployable Solar Panel, Sensorarray,  MDDE (+2), Base AI system, Mech sized Weaponry (Advanced
    (Vibro/electro/Monofilament) Sword), Rocketlauncher, Water venting system, midtech Heatsink}\\&&\\
    H    &    7        &    \makecell*[{{p{12cm}}}]{BCI, Advanced AI, WeaponLaser, Sustained Flight,  Advanced
    Shields, Nano Repair Cloud, Automedic, OVERDRIVE mode, Hightech Heatsink}\\&&\\
    E    &    8        &    \makecell*[{{p{12cm}}}]{Quantum Lookahead Sensors, Blink Teleporter, Timerift Shields
    (Delays Damage 1d10 rounds), EnergySword}
\end{tabular}\par
Designtarget (TODO: write words good) for normal fighting mech

\begin{enumerate}[label = - ]
\item 20\% movement
\item 20\% weaponry
\item 5\% energy
\item 5\% Cargo
\item 50\% + armor
\end{enumerate}


\subsection{Energy Systems}\label{subsec:energy-systems}
Energy Systems usually do not use a hardpoint.

\textbf{Power to Weight} is how much maximum output power per weight the system provides.

\textbf{Input} of an Energy System is generally considered its "fuel", but can also just be batteries or similar.
The Input is stored with other Cargo.

\textbf{Efficiency} is the function with input x being ratio of poweroutput vs maximum power output and result of
efficiency.
In parenthesis at the end is the efficiency at 100\%

\textbf{Output} is set during addition/modification of the system
(Engineering:Repair+ system specific skill, Competence vs 8, 4h).
By Default there are only 2 Modes: Power saving and Full Throttle.
Mode usage is marked for full hours spent in them, with the higher usage overriding the lower usage, except when the
remaining Energy is less than an hour, where minute/round calculations might be needed.
Energy E is in Rates of kilojoules/round, since Rounds are 5 seconds long, 1 Energy is equivalent to 5kw over 1 Round.
1 MJ is 1000 kJ so 1 MJ = 200E

\begin{tabular}{l|cccll}
    Type & E per t & heat/r per t & input & efficiency\\
    \hline Basetech \\
    Steampower & 0.7 & 8 & Coal & 10\%\\
    FlyWheel & 20 & 0 & stored & loses 10\% per turn \\
    Crew & 0.01 & 0 & Labor & 10\%\\
    \hline Lowtech \\
    Combustion Motor& 140 & 2  & Combustible & 40\% \\
    \hline Midtech \\
    PEM & 100 & 0.1 & Hydrogen & 90\%\\
    Solar Panel & 10 & 2 & Sunlight & Contamination dependent\\
    Batteries & 10 & 2 & battery charge & loses 10\% per month \\
    \hline Hightech \\
    Reactor & 300 & 5  & Fuelrods & 8\%\\
    CFReactor & 40 & 0.1  & Deuterium & 80\%\\
    Capacitor & 1000 & X & electricity & loses 10\% per hour
\end{tabular}


\textbf {Solar Panel} Deploying takes 5 turns times 100w or 5 minutes labor, any damage received while deployed
will also damage the Solar panel, production halved every 5 contamination levels, only works during day

\textbf {CFReactor} Cold Fusion requires start up energy, equaling the Output for 1 Round.

\textbf {Deuterium} 1 kg - 4.5 TJ or 900,000,000E (effective at 80\%= 720ME)

\textbf {Fuel Rod} 1 kg - 75 TJ or 15,000,000,000E (effective at 8\%= 120ME)

\textbf {Hydrogen} 70kg (1000L) = 10,082MJ or 2,016,400E (effective at 90\% = 1,814,760E)

\textbf {HFC-10} 10 liter - hydrogen fuel cell: Cylinder 20cm diameter,
35 cm height or 23cm cube=10L. \par
containing 9 MJ or 1800E

\textbf {HFC-50} 50 liter - hydrogen fuel cell: Cylinder 40cm diameter,
40 cm height or 37cm cube = 50L,
\par containing 45 MJ or 9000E\par
contains low power sub system used to power up basic systems (like hydrogen gas pumps for the PEM Stacks)

\textbf{CombustionFuel} 800kg (1000L) = 40,000MJ or 8,000,000E (3,2ME)

\textbf{Battery} has 180kE/ t

\textbf{Capacitor} generates 1 Heat per 10 Energy moved, for all other purposes instant

\subsection{Movement Systems}\label{subsec:movementsystems}

\textbf{Efficiency} is the ratio of power consumed to power used to move the mech

\textbf{Drag} is the general amount of power wasted the faster the mech is going

\textbf{Power to Weight} is how much maximum movement power the system can employ per weight

\begin{tabular}{l|ccccccl}
    Type    & Energy/r & GroundCoefficient per ton & Area per sector & Power per ton & hardpoints &Extra\\
    \hline Basetech\\
    railbased            & 55 & 0.001 & 20 & 500 & 50\% & rails \\
    simple wheels        & 100 & 8 & 1 & 750 & 25\% & flat terrain \\
    floatation system    & 15 & 2 & 1000 & 90 & 10\% & water\\
    \hline Lowtech\\
    complex wheels       & 175 & 6 & 3 & 1500 & 25\% &  somewhatflat terrain \\
    tracks               & 100 & 0.02 & 40 & 600 & 33\% & rough terrain\\
    2legs                & 30 & 0.5 & 100 & 450 & 25\% & all terrain\\
    4legs                & 30 & 0.2 & 200 & 450 & 25\% & all terrain\\
    6legs                & 30 & 0.13 & 300 & 450 & 25\% & all terrain\\
    jumpjets             & 1000 & 50 & 1 & 1000 & 1 & upwards\\
    \hline Midtech\\
    2 flexible legs        & 100 & 10 & 10 & 1000 & 30\% &  like a large human\\
    4 flexible legs        & 100 & 1 & 100 & 1000 & 30\% &  like a large human\\
    6 flexible legs        & 100 & 0.3 & 1000 & 1000 & 30\% &  like a large human\\
    boostjets            & 2000 & 5 & 0.1 & 1000 & 2 &  any direction\\
    \hline Hightech\\
    Spiderwalker         & 150 & 0.1 & 500 & 400 & 50\% & any surface
\end{tabular}

Accelleration is assumed to be able to change to 50\%/80\%/100\% of max speed  in the first/second/third round\par
Pneumatic variants of the legs weigh half as much but are constrained to weightclasses light and below\par

\subsubsection{Maximum Speed}
Since Power is in Watt, aka
\(\frac{kg*m*m}{s*s*s}\)
``all'' that is needed to get the max speed in
\(\frac{m}{s}\) is to get
the friction forces at the point of highest velocity in \(\frac{kg*m}{s*s}\), aka Newton.
The GroundCoefficient is the one given in the Table, Power has to be calculated
\(Friction = Mass* GroundCoefficient/100 * 9.81\)\par
Repeat updating Airdrag\par
\(Airdrag = \frac{Airdrag + 0.5*Area*1.225 * Speed^2}{2}\)
and Speed\par
\(Speed = \frac{Power}{Airdrag + Friction}\)\par
Until the maximum value is found.\par
Bonus Formulae:\par
\(Airdrag = \frac(49*Speed^2*Area)/80 \) \par
\(Speed = \frac{Power}{(49*Speed^2*Area)/80 + Friction}\)


\subsection{Heat Systems}\label{subsec:heat-systems}
Warning: ALL Dissipated heat contributes to IR signature
\textbf{Input} is the amount of power the Heat System requires to function actively
\textbf{Weight} direct weight of the system
\textbf{Cooling} how much heat is dissipated, passively/actively.
\textbf{Capacity} how much heat can be stored in this system.
Cumulative for the mech.

\begin{tabular}{r|lllll}
    Type & Input & Weight & Cooling & Capacity &Extra\\
    \hline Basetech\\
    Small Heatsink & 0 & 0.5t & 0.2 & 15 & Storage\\
    \hline Lowtech\\
    Radiator Vent & 2 E/r &scalable &(1/2) + (1/t / 5/t) & 0 & 1 Hardpoint\\
    Large Heatsink & 0 & 1t & 0.5 & 50 &Storage\\
    \hline Midtech\\
    Optimized heatsink  & 0 & 1t & 1 & 80 &Storage\\
    Coolant Dump  & 200 E & 4t & 0/100 & 100 & Decoy; 1 Hardpoint \\
    Radiator Wings & 0.2 E/r & 2t & 10\% & 100 & Deploy; 1 Hardpoint\\
    \hline Hightech\\
    Fractal heatsink & 20 E/r & 2t & 3/0 & 200 & Storage\\
    OVERDRIVE & 1000E & 1t & 0 & & Gadget
\end{tabular}\par

\textbf{Coolant Dump}: When triggered, releases 100 stored heat together with 100L water, as a thermal barrier

\textbf{Radiator Wings}: Deploying takes 5 turns and 1E or 5 minutes labor, any damage received while deployed
will also damage the Radiator Wings: radiates 10\% of total system heat, as long as 0.2E/r (for coolant pumping)
is applied.

\textbf{Overdrive}: When triggered, enter overdrive mode: able to act at double speed,
stored heat is spent insteadof gained.
When heat drops to 0 or below, Mech shuts down for 2d10 rounds and Overdrive requires a Competence, Gadgets-Repair;
10m;
Threshhold of 10 to repair.

\subsection{Sealsystems}\label{subsec:sealsystems}
\textbf{Level} how much the seal reduces Contamination behind it.
\textbf{Resistance} how much damage it takes to damage the Seal.
For each Instance of Damage at or above the Resistance, it's Resistance is lowered by 1.

\begin{tabular}{l|r|r|r}
    Tech & Seal Level & Resistance & Cost/Sector\\\hline
    Base &  6 & 2 & 1\\
    Low &   13& 3 & 2\\
    Mid &   21& 4 & 3\\
    High &  30& 5 & 4\\
    Experimental& 50 & 2 &8
\end{tabular}\par

\subsection{Weapons}\label{subsec:weapons}
\textbf{Damage}
\textbf{Threshholds}
\textbf{Costs}
\textbf{Skills}

Melee weapons rely on the movement System, so their costs are in meters that are used to attack rather than move\par
\begin{tabular}{c|clll}
    Type & Damage & Threshholds & Costs & Weight\\
    \hline \hline BaseTech&&&&\\
    \hline Metal Sword & 10 & 10m[5,7,9,11,13,15,17,19,20]0 & 30m & 5t \\
    \hline Manned Turret & X & X & \makecell[l]{manual \\labor +X} &  X +50\% \\
    \hline \hline Lowtech &&&&\\
    \hline Heavy Machinegun & 30 & \makecell[l]{100m[3,6,10,14,18]20\\
                                        500m[5, 8, 12,15]10\\
                                        1000m[7,10,14]5\\
                                        2km[11,16]5\\
                                        8km[15]1 indirect}
                                            & 5kg, 5H & 100kg \\
    \hline Light Artillery & 120 & \makecell[l]{7km[8,12,16,20]1 \\
                                        15km[10,14]}& 50kg, 10H & 8t  \\
    \hline 120mm Smoothbore & 120 & \makecell[l]{50m[10,20]3\\
                                        100m[6,10,14]8\\
                                        500m[5,7,9,11]15\\
                                        2km[8,10,12,14,16]20\\
                                        4km[11,13,15,20]10 indirect\\
                                        8km[13,17]5 indirect}
                                            & 20kg 5H & 20t \\
    \hline\hline Midtech &&&&\\
    \hline Autocannon & 80 & \makecell[l]{10m[10]3\\
                                  200m[5,8,11,14,17,20]15\\
                                  1km[8,12,16,20]20\\
                                  2km[10,14,18]10 \\
                                  4km[13,17,20]5 indirect }
                                        & 7kg 5H & 5t \\
    \hline Pulse Laser & 60 & \makecell[l]{1km[5,-15]50\\
                                   5km[7,9,11,13,15]100\\
                                   LOS[10,12,14,16,18]100}
                                        & 1000E 10H & 7t\\
    \hline Precision Cannon & 140 & \makecell[l]{500m[7,10,13]2\\
                                         2km[7-11,13,15,17,19]30\\
                                         LOS[8,10,12,14,16,18,20]25}
                                        & 15 kg 10H & 10t\\
    \hline Rocket Launcher x 8 & X & X & X & 500kg \\
    \hline\hline Hightech &&&&\\
    \hline Rocket Platform & X & X & X & 1500 kg \\
    \hline Railgun & 160 & \makecell[l]{200m[7,9,11,13,15]5\\
                                1km[5,7,9,11,13]10\\
                                4km[7,9,11,13,15]5\\
                                12km[9,11,13,15,17]1}
                                & 50 kg, 25H, 4000E & 15t\\
    \hline Advanced Pulse Laser & 90 & \makecell[l]{3km[5-20]50\\
                                   9km[7,9,11,13,15,17]100\\
                                   LOS[10,12,14,16,18]100}
                                        & 1800E 15H & 3t\\
    \hline Ramping Beam Laser & 10\(\times\)X & \makecell[l]{1km[5,10,15]10\\
                                          LOS[10,15]100}
                        & X\(\times\)(10E 1H) & 10t\\

\end{tabular}\par

\textbf{Notes}:
\begin{enumerate}[label = - ]
\item Melee Weapons: Ignore Shields
\item Manned Turret: X Depends on the mounted gun and the skill of the gunner
\item Light Artillery: Attacks take 1s per 300m distance, rounded up, minimum range 100m
\item 120mm Smoothbore: Can only fire every second round, requires Indirect for targets further away than 2k
\item Autocannon: may make another attack instantly for double the cost
\item Basic Pulse Laser: -0.1*contamination level dmg per km, rounded mathematically
\item Precision Cannon: Needs to be stationary to aim and fire, 1 turn of aiming between shots
\item Rocket Launcher: Ignores Shields
\item Rocket Platform: Ignores Shields, Beam rules
\item Railgun: Needs to be stationary to fire.
Loading takes 2 turns.
\item Advanced Pulse Laser: -0.1*contamination level dmg per km, rounded mathematically
\item Ramping Beam Laser: Every round, X increases by 1 while shooting, decreases by 1 down to 0 otherwise.
\end{enumerate}


\subsection{Defensesystems}\label{subsec:defensesystems}
Defensesystems are configured using the relevant \hyperref[subsubsec:defense]{Technical Skill},
the relevant usage Method, usually \hyperref[subsec:computer]{Computer Usage}
with (if not explicitly defined) the threshholds [5,8,13,16,19] or automatically by a Program of the corresponding level.


Configuration, Strength are the  armor-equivalent statistics a shield provides.
Configuration is required to start a Shield and takes 1 Action, if the lowest threshhold is not hit, the Configuration
is 0, but the Shield still starts.

Coldboot can only be done after all Systems have been shut off for at least one turn.

Shields turn off when overwhelmed and have to meet the Reboot requirements (including another Configuration) to turn
back on.
Config X means that the specific Configuration Threshhold of X has to be reached.
Each Sector requires a Defensesystem on its own.
Defensesystems provide coverage for the Sector that they are in first, but if multiple levels (representing a larger
Generator) are present they can (during configuration) be spread to adjacent Sectors.
Deflectors and Dampeners provide their effects, if they do not cover all sectors, by rounding the portion of sectors
that are uncovered up to the nearest tenth and rolling a single die (Fate point adds Fate Bonusdice).
If the result is larger than the uncovered portion, the Field generator applies.
(1/8 missing requires 3+, 7/8 missing requires 10)

\begin{tabular}{c|cccccc}
    Type & Protection &Weight& Cost & Failure  & Reboot & Coldboot \\
    \hline Basetech\\
    Applique-Armor & Strength 1 & 1t & - & destroy, block & - & permanent \\
    \hline Lowtech\\
    Basic Deflector & 1 Evasion & 0.2t & 10E/r, 1H/attack & - & - & 5E, 1r, Config 5\\
    Basic Dampener & 5 + Config Reduction & 0.5t & 20E/r, 3H/attack & - & - & 30E, 3r \\
    \hline Midtech\\
    Shield & Strength 3 & 1t & 40E/r or 10E/r, 1H/attack & passthrough & 5r & 2000E \\
    Complex Deflector & 3 Evasion & 0.8t & 40E/r, 2H/attack & - & - & 160E, 8r, Config 10\\
    Heavy Dampener & 30 + Config\(\times 3\) Reduction & 1t & 60E/r, 5H/attack & - & - & 100E, 10r\\
    Flash Shield & Strength 1 & 1t & 20E/r), 2H/attack & block & X\(\times\)50E & 200E, 2r\\
    \hline Hightech\\
    Advanced Shield & Strength 4 & 1t & 200E/r or 20E/r & block & 1r & 5r
\end{tabular}

\vspace{1cm}

\textbf{Deflectors} need to be Configured with the Deflector skill and require a minimum degree of success to
work and provide flat evasion-level bonuses.

\textbf{Dampeners} need to be Configured with the Shield skill and provide a flat damage reduction to all
attacks that are influenced by shields, they usually scale with the Quality of the Configuration.

\textbf{Flash Shield} As soon as the Reboot Energy is paid and the configuration is done, they turn back on.
X is set to 0 at Coldboot and increased by 1, everytime the Shield is broken.
A 10-Resonance on the Configuration reduces X by that much, even going negative.
While X is negative, the Energycost is 0.

Cost: per sector for being active.

Failure: Normal behaviour on Failure is for the shield to turn off until rebooted.

Dampeners and Deflectors dont fail from damage.

\subsection{Armor}\label{subsec:armor}
\begin{tabular}{c|ccc}
    Type & Strength & Weight & Failure\\
    \hline BaseTech&&&\\
    Outside Cargo & 0 & 1.1t & destroyed, passthrough\\
    Metal Plating & 2 & 2t & damaged, passthrough\\
    Ultra High Density Composite & 4 & 5t & out of order, block\\
    \hline Lowtech &&&\\
    Composite & 3 & 2.5t & damaged, passthrough \\
    \hline Midtech &&&\\
    Active & 1/4 & 2t & out of order, passthrough full/1\\
    SpeedShell & 1 & 0.5t & destroyed, passthrough \\
    \hline Hightech &&&\\
    Ion & 5 & 1t & destroyed, energy system damage, block 
\end{tabular}

\begin{enumerate}[label = - ]
    \item Outside Cargo: Provides flat reduction depending on cargo.
    Cargo gets destroyed.
    \item Metal Plating: Repairs need 12 or more to not lower Quality
    \item Ultra High Density Composite: Repairs need 8 to succeed and Freq1 Resonance lowers Strength by Amp
    \item Composite: needs to be replaced, spare parts 200kg, success with 6+, scrap at 2
    \item Active: needs 2E per round and per incoming hit or Strength is 1.
    Repairs Success with 15+/Spare 5+, 1--4 Scraps the Layer, Freq1 subtracts
    \item Speedshell: Vent 1H per 5m/s;
    Repair 12+, 1--11 Scraps the Layer
    \item Ion: 4E/r per Segment, on destruction energy Output is reduced by 10\% until scrapped armor is removed.
\end{enumerate}



\section{Misc Modules}\label{sec:miscmodules}
\par
\subsection{Basetech}\label{subsec:basetech}
\subsubsection{Water Distillery}\label{subsubsec:water-distillery}
Energy Cost: 50kw\par
Heat: 10\par
Weight: 0.5t\par
Turns 1L of Dirty Water into 1L of Contaminated Water per minute\par
Requires Reprocessing\par
\par

\subsubsection{Cargo compartment}\label{subsubsec:cargo-compartment}
stores up to 1 t of Cargo\par
Weight: 0.5t for Size Calculation, Empty weight 0.1t, Full 1.1t for speed calculation (if relevantt)
\par
\section{Lowtech}\label{sec:lowtech}
\subsubsection{Basic Food Processor}
Energy Cost per operation: 10kw \par
Heat: 5 \par
Size:0.5t\par
Detoxes and prepares up to 5kg of Food in 30 minutes \par
Requires Reprocessing\par

\subsubsection{Basic Air Filter}
Energy Cost per operation: 15kw \par
Heat: 2\par
Size:0.5\par
Decontaminates and pressurizes 3000 L of air in  10 minutes or double without pressurizing. \par
Requires Decon 1\par
Airtank 10 Liter @ 300 bar contains 3000L of air.
Air is used at 10 Liters/Minute so a 10 Liter air tank lasts about 5 hours of calm
walking or half (or even less) when doing strenuous activity.
1.225 kg/m3: 3.675 Kg of air
+ 15kg for Tank, Valve und Hose
\par
\subsubsection{Basic Decon Equipment}
Energy Cost per operation: 20 kw \par
Clean Water Cost 10:1\par
Heat: 10 \par
Size: 0.5\par
Allows to lower Objects contamination by washing them, if applicable.
Rate requires 1 hour per Kilogram, Contamination is 1/10 of starting value.
Used water is Contaminated\par

Optical Sensors\par
Energy: 1\par
Size: 0.25 \par
Rating: 100\par
Provides Targetting Lock is obtained\par
Visual Lock: Distance/Rating <= Target Size (+Camouflage and visual impairments)\par

Thermal Sensors\par
Energy: 1\par
Size 0.25\par
Rating 100\par
Thermal Lock: Distance/Rating <= Target Thermal Signature\par

\subsubsection{Basic Computer System}\par
Weight 500kg\par
Energy: 0.1\par
Module: Processor or Memory\par
\par
\section{Midtech}\label{sec:midtech}
\subsubsection{ComputerSystem}:
Weight 100 kg
Energy: 0.1
Module: Processor or Memory\par

Nondeterministic Computation Core:\par
Size 0.5\par
Energy: 0.5\par
Module: Processor\par
Requires Black Biotech 3\par

Extended Decon Equipment\par
Energy Cost per Operation: 36k (10/s)\par
Alacast Cost per Operation: 5g\par
Size: 1
Allows to turn Contaminated Objects into uncontaminated Objects at the rate of 100kg/h if suitable\par
If used on Humans, Roll Red Biotech, on 9 or more, removes one Charactertaint.\par
Always Inflicts Wound of Level 11-RollResult\par
Requires Decon 3\par

Extended Detox Equipment:\par
Energy Cost per Operation: 36k (10/s)\par
Size 0.75\par
Allows Detoxing and Preparing of up to 1t/h\par
If used on Humans, Roll Red Biotech, on 6 or more, remove all Infections and appropriate Toxins.\par
Always Inflicts Wound of Level 8-RollResult\par
Requires Detox 3\par
Automedic:
Energy Cost: 6k(variable - 10/s for 10 min)\par
Size: 0.75\par
Supplies tools of Medic/ supplies good operating conditions for Medic / reduces the use of Medicine by half / allows stasis\par
Roll Red Biotech, if the roll surpasses the wound, supplies an operating character with a temporary Medic skill and
double Advantage or autonomously attempts the procedure at Level 3 mod -1.\par
If not, Instead supplies one Advantage.\par
Software may supply the Red Biotech roll even autonomously.\par
Stasis is automatic and slows down the daily/hourly tickrate (including Healing and Contamination) by a factor of 10
(additive) for every energy point per second.
This is a medical coma, not timedilation.
\section{Hightech}\label{sec:hightech}
Advanced Detox-Decon Equipment \par
Size: 1\par
Energy Cost per Operation 72k (variable 20/s for 1 h)\par
Alacast Usage per Operation if Detainting: 1g/t of raw material, 10g/human\par
Detoxes and Detaints up to 5t of material, including air and water.\par
Living Entities suffer no harm, and on a 4 or more on a Red Biotech Roll lose one Taint level.
Furthermore all Infections and Appropriate Toxins are removed.\par
Requires Detox 4 and Decon 4

Advanced AutoMedic\par
Energy Cost: 3k (variable - 5/s for 10 min)\par
size: 1

Allows Medic Rolls/ supplies very good operating conditions for Medic Rolls/ reduces the use of Medicine by 90\%/allows Ressurection\par
To Treat a wound:\par
Roll Red Biotech, if the roll surpasses the wound, supplies an operating character with a temporary Medic skill and
triple Advantage or autonomously attempts the procedure at Level 4 mod -1.\par
If not, Instead supplies one Advantage or autonomously attempts the procedure at Level 2.\par
Software may supply the Red Biotech roll even autonomously.\par
To Ressurect:\par
requires 1g of EOL per total spent points on a character (or 20* the average level of an NPC)
Roll Black Biotech, every 15\% missing from the corpse or order of magnitude in days that has passed since the time of death
lowers the final score by 1.
If the roll succeeds, place 5 wounds of severity 15-result on that character, they lose
one level from the highest Ability and one from the highest Skill and are alive again.
If the roll does not succeed,
5\% of the total corpse and the dosage of EOL is lost.


\end{document}
